%!TEX TS-program = xelatex
\documentclass{abgabeOTH}
\usepackage[utf8]{inputenc}
\usepackage[german]{babel}
\usepackage[doublespacing]{setspace}
\usepackage{tabularx}
\usepackage{stix}
\usepackage{amssymb}
%\usepackage{ragged2e}
\makeatother
%-----Allgemeine Variablen----%

%Masterarbeit oder Bachelorarbeit
\newcommand\MaBaArbeit{Bachelorarbeit}
%Name des Studenten Eintragen(Name, Vorname)
\newcommand\Student{Superfleissiger Student}
%Studiengang des Studenten eintragen
\newcommand\Studiengang{Industrie 4.0 - Informatik}
%Matrikelnummer
\newcommand\Matrikelnummer{123445678}
%Semester eintragen
\newcommand\Semester{8}
%Ausgabedatum der Arbeit
\newcommand\Ausgabe{01.01.2020}
%Festgelegetes Abgabedatum
\newcommand\Abgabe{01.01.2025}
%Titel der Arbeit eintragen
\newcommand\Titel{Megatoller Titel in deutsch}
\newcommand\TitelEnglisch{Super awsome title in english}
%Kurzbeschreibung
\newcommand\Beschreibung{Total tolle und aussagekräftige Beschreibung der megaspannenden Aufgabenstellung der Abschlussarbeit mit dem super genialen Titel}
%Name des Erstprüfers
\newcommand\Aufgabensteller{Prof. Dr. Schlaubi Schlumpf}
%Name des Zweitprüfers
\newcommand\Zweitpruefer{Prof. Dr. Gargamel}



%----externe Bachelorarbeit----%
%Name der Firma
\newcommand\Firma{FIRMA}
%Firmenanschrift
\newcommand\Anschrift{ANSCHRIFT}
%Name des Betreuers innerhalb der Firma
\newcommand\Firmenbetreuer{NAME}
%Abteilung der Firma
\newcommand\Abteilung{ABTEILUNG}
%Telefonnummer
\newcommand\Telefon{12345678}



%--Änderungen--
%Zeitraum um den die Abgabe verlängert werden soll in Zahlen
\newcommand\Zeitraum{Wochen}
%Neuer Abgabetermin
\newcommand\AbgabeNeu{01.01.3000}
%Gründe für die Verzögerung
\newcommand\Grund{akzeptable Gründe sind z.B.: unvorhergesehene Komplexität des Themas,schwierigkeiten bei der Recherche(keine Literatur vorhanden), lange nachgewiesene Krankheit, Pflege von Familienangehörigen(nicht Kinderbetreuung)}





\makeatletter
%Um das entsprechende Feld anzukreuzen einfach \mybox durch \myboxcheck ersetzen
\newcommand{\mybox}{$\square$}
\newcommand{\myboxcheck}{${\Large\rlap{ \checkmark}}\square$}

\begin{document}

\vspace*{.3cm}
\subsection*{Annahmebestätigung der \MaBaArbeit}
\vspace{-.3cm}
\noindent\rule[1ex]{\textwidth}{.5pt} %für durchgezogene Linie
    An das Dekanat (4 Exemplare, in Maschinenschrift ausgefüllt) \hfill Eingang:\rule{3cm}{.5pt}\\
    Ich habe heute von der Studentin/ dem Studenten\\
    (Name, Vorname)\hfill \Student\\
    Matrikelnummer \hfill \Matrikelnummer\\
    Studiengang:\hfill \Studiengang\\
    \raggedright Bei Erfüllung der Voraussetzung möchte ich folgende Vertiefungsrichtung im Zeugnis eingetragen\\
    haben(trifft nur auf EI-Studierende zu):\\
    \centering{{\Large \mybox} \small{AUT} \qquad  {\Large \mybox} \small{CPS} \qquad {\Large \mybox} \small{ENT} \qquad {\Large \mybox} \small{IKT}} \\
    \raggedright
    Semester:\quad \Semester\\
    eine Bachelorarbeit mit dem Titel\\
    Deutscher Titel:\quad \Titel\\
    Englischer Titel:\quad \TitelEnglisch\\
    angenommen.\\
    Ausgabetag war:\quad \Ausgabe\\
    Letzter Abgabetermin:\quad \Abgabe\\
    Ich bestätige die Annahme des praktischen Teils der Bachelorarbeit:\hfill Ja \mybox \qquad Nein \mybox \\
    \raggedleft Nicht erforderlich \mybox \\
    öffentliche Auslage:\hfill  Ja \mybox \qquad Nein \mybox\\
    falls Nein: nicht öffentliche Auslage bei:\hfill \Aufgabensteller \\
    \hfill Zeitraum:\rule{3cm}{0pt}\\
    Aufgabensteller:\hfill \Aufgabensteller\\
    Aufgabensteller (Unterschrift): \hfill \rule{.3\textwidth}{.5pt}\\
    Bestätigung(Unterschrift) der Studentin/ des Studenten\hfill \rule{.3\textwidth}{.5pt}
    \vspace*{.5cm}
    \raggedright

   % \rule{\textwidth}{.5pt}
   % \tiny{Verteiler: Dekanat \qquad Aufgabensteller \qquad Studentin/Student \qquad Studentenamt}

\end{document}
